\documentclass{article}
\usepackage[preprint]{neurips_2025}
\usepackage[utf8]{inputenc}
\usepackage[T1]{fontenc}
\usepackage{hyperref}
\usepackage{url}
\usepackage{booktabs}
\usepackage{amsfonts}
\usepackage{amsmath}
\usepackage{amssymb}
\usepackage{graphicx}
\usepackage{xcolor}
\usepackage{multirow}
\usepackage{natbib}
\usepackage{microtype}

\newcommand{\speccheck}{\textsc{SpecCheck}}
\newcommand{\aucroc}{\text{AUC-ROC}}
\newcommand{\aucpr}{\text{AUC-PR}}

\title{SpecCheck: Testing Confidence Monotonicity Across Specificity Levels\\for LLM Hallucination Detection---A Negative Result}

\author{
  Anonymous Author(s)
}

\begin{document}

\maketitle

\begin{abstract}
Large language models (LLMs) frequently hallucinate specific details while retaining approximately correct general knowledge, suggesting that a model's confidence should increase monotonically as a claim is made more abstract.
We propose \speccheck{}, a training-free framework that tests this \emph{confidence monotonicity} hypothesis by generating a ``specificity ladder'' for each atomic claim---a sequence of progressively abstracted versions---and checking whether the model's logprob-based confidence increases monotonically.
We evaluate \speccheck{} on three benchmarks (FActScore, LongFact, TruthfulQA) across three open-source LLMs (Llama-3.1-8B, Mistral-7B, Qwen2.5-7B).
Our results are largely negative: \speccheck{} achieves AUC-ROC of 0.48--0.62 across settings, consistently underperforming SelfCheckGPT (0.51--0.73) and logprob confidence baselines.
The confidence monotonicity hypothesis holds for the majority of claims regardless of factuality, yielding insufficient discriminative power.
However, we find that \speccheck{} provides a modest complementary signal when combined with existing methods via logistic regression (+0.3--2.1 AUC-ROC points on FActScore).
We present this work as an honest analysis of a principled but unsuccessful approach, providing insights into why specificity-based probing fails and when it may offer marginal value.
\end{abstract}

%% ============================================================
\section{Introduction}
\label{sec:intro}
%% ============================================================

Large language models generate fluent text that frequently mixes accurate information with fabricated details---so-called hallucinations~\citep{min2023factscore, manakul2023selfcheckgpt}.
A common pattern is that models hallucinate at the level of \emph{specific details} (exact dates, precise numbers, particular names) while retaining approximately correct general knowledge (the rough time period, the order of magnitude, the category of entity).
For example, a model might correctly associate Einstein with the Nobel Prize in the early 20th century but hallucinate the exact year or the specific field.

This observation motivates a natural hypothesis: if a model truly knows a fact, its confidence should be monotonically non-decreasing as the claim becomes more abstract.
A model should be at least as confident about ``Einstein won the Nobel Prize in the 1920s'' as about ``Einstein won the Nobel Prize in 1921.''
Violations of this \emph{confidence monotonicity} across a \emph{specificity ladder}---a sequence of the same claim at decreasing levels of detail---could serve as a hallucination signal.

Existing hallucination detection methods operate at a single level of specificity: they verify claims against external knowledge~\citep{min2023factscore, wei2024longform} or check cross-sample consistency~\citep{manakul2023selfcheckgpt, farquhar2024detecting}.
None exploit the specificity structure of claims.

We propose \speccheck{}, a training-free framework that (1)~decomposes LLM outputs into atomic claims, (2)~generates a specificity ladder for each claim via prompted abstraction, (3)~estimates confidence at each level via logprob-based yes/no probing, and (4)~detects hallucinations by testing for monotonicity violations.

\textbf{Our results are largely negative.}
\speccheck{} underperforms SelfCheckGPT and simple logprob baselines across all three benchmarks and all three models.
The core issue is that confidence monotonicity holds for the majority of claims \emph{regardless of factuality}---both true and hallucinated claims exhibit similar monotonicity patterns, providing insufficient discriminative power.

Our contributions are:
\begin{itemize}
    \item We propose and rigorously test the confidence monotonicity hypothesis for hallucination detection, providing a principled negative result.
    \item We conduct extensive experiments across 3~models, 3~benchmarks, and multiple ablations, demonstrating where and why the approach fails.
    \item We show that \speccheck{} provides a modest complementary signal when combined with existing methods, suggesting the specificity dimension captures some independent information.
\end{itemize}

The remainder of this paper is organized as follows.
Section~\ref{sec:related} reviews related work.
Section~\ref{sec:method} describes the \speccheck{} framework.
Section~\ref{sec:experiments} presents our experimental setup and results.
Section~\ref{sec:analysis} provides analysis and ablations.
Section~\ref{sec:discussion} discusses limitations.
Section~\ref{sec:conclusion} concludes.


%% ============================================================
\section{Related Work}
\label{sec:related}
%% ============================================================

\paragraph{Hallucination Detection.}
FActScore~\citep{min2023factscore} pioneered atomic claim decomposition and verification against a knowledge source.
SAFE~\citep{wei2024longform} extends this with search-augmented verification.
Both require external knowledge bases, whereas \speccheck{} is self-contained.
SelfCheckGPT~\citep{manakul2023selfcheckgpt} detects hallucinations by checking consistency across multiple sampled outputs---if a fact appears consistently across independent samples, it is likely known by the model.
Semantic entropy~\citep{farquhar2024detecting} clusters semantically equivalent responses and computes entropy over meaning clusters, providing answer-level uncertainty.
\speccheck{} differs by probing the \emph{same claim at different specificity levels} rather than comparing different complete responses or relying on external verification.

\paragraph{Uncertainty and Calibration.}
Verbalized confidence methods elicit numerical confidence scores directly from LLMs~\citep{kadavath2022language}.
Self-consistency~\citep{wang2023selfconsistency} uses sampling agreement as a proxy for confidence.
These operate at the response level; \speccheck{} extends self-assessment to the claim level across multiple granularities.

\paragraph{Abstraction in LLMs.}
Step-Back Prompting~\citep{zheng2024stepback} uses abstraction to improve reasoning by first answering a more general question.
\speccheck{} applies a similar abstraction idea but for \emph{detection} rather than generation, and does so systematically across multiple levels.

Unlike prior work, \speccheck{} is the first method to probe model confidence across a spectrum of specificity levels for the same underlying fact.
Our negative results reveal important limitations of this approach.


%% ============================================================
\section{Method}
\label{sec:method}
%% ============================================================

\speccheck{} operates in four stages.
Given an LLM-generated passage, we:
(1)~decompose it into atomic claims,
(2)~generate a specificity ladder for each claim,
(3)~estimate confidence at each specificity level, and
(4)~detect hallucinations via monotonicity testing.

\subsection{Claim Decomposition}

Following FActScore~\citep{min2023factscore}, we decompose long-form generations into atomic claims using a structured prompt.
Each atomic claim is a single, self-contained factual assertion (e.g., ``Marie Curie won the Nobel Prize in Physics in 1903'').

\subsection{Specificity Ladder Generation}

For each atomic claim $c_0$ (the original, most specific version), we prompt the same LLM to generate a sequence $c_1, c_2, c_3$ at decreasing specificity levels:
\begin{itemize}
    \item \textbf{Level 0 (Original)}: The exact claim. E.g., ``Marie Curie won the Nobel Prize in Physics in 1903.''
    \item \textbf{Level 1 (Approximate)}: Precise values replaced with ranges. E.g., ``Marie Curie won the Nobel Prize in Physics in the early 1900s.''
    \item \textbf{Level 2 (Category)}: Specific instances replaced with categories. E.g., ``Marie Curie won a Nobel Prize in a science field.''
    \item \textbf{Level 3 (Abstract)}: Most general version. E.g., ``Marie Curie won a major scientific award.''
\end{itemize}

We use few-shot prompting with diverse examples covering numerical, temporal, entity, and categorical claim types.
For claims with clear numerical values, we also employ template-based abstraction as a fallback.
Generation uses temperature~0 for deterministic ladder construction.

\subsection{Multi-Granularity Confidence Estimation}

For each claim version $c_k$ at level $k \in \{0, 1, 2, 3\}$, we estimate the model's confidence via logprob-based probing.
Specifically, we formulate a yes/no question: ``Is the following statement true? \texttt{"}$c_k$\texttt{"} Answer with Yes or No.''
We run a single forward pass and extract the softmax probability over the ``Yes'' and ``No'' tokens:
\begin{equation}
    \text{conf}(c_k) = \frac{\exp(\text{logit}_{\text{Yes}})}{\exp(\text{logit}_{\text{Yes}}) + \exp(\text{logit}_{\text{No}})}
\end{equation}
We sum probabilities across token variants (``Yes''/``yes''/`` Yes'').

\subsection{Monotonicity-Based Detection}

Given the confidence sequence $[\text{conf}(c_0), \text{conf}(c_1), \text{conf}(c_2), \text{conf}(c_3)]$, we define:
\begin{align}
    M(c_0) &= \frac{1}{3} \sum_{k=1}^{3} \mathbf{1}[\text{conf}(c_k) \geq \text{conf}(c_{k-1})] \label{eq:mono} \\
    S(c_0) &= 1 - M(c_0) + \alpha \cdot (\text{conf}(c_3) - \text{conf}(c_0)) \label{eq:speccheck}
\end{align}
where $M$ is the monotonicity score (fraction of adjacent pairs satisfying monotonicity), $S$ is the \speccheck{} score combining monotonicity violations with the confidence gap, and $\alpha = 0.5$ is a weighting parameter.
A high \speccheck{} score indicates potential hallucination.

\textbf{Hypothesis.} If a model truly knows a fact, confidence should be monotonically non-decreasing as specificity decreases (i.e., $M = 1$).
Violations ($M < 1$) signal that the model's knowledge is unreliable at the specific level of detail.


%% ============================================================
\section{Experiments}
\label{sec:experiments}
%% ============================================================

\subsection{Setup}

\paragraph{Models.}
We evaluate three open-source instruction-tuned LLMs: Llama-3.1-8B-Instruct, Mistral-7B-Instruct-v0.3, and Qwen2.5-7B-Instruct, representing three major model families.
All models run in float16 on a single NVIDIA A6000 (48GB).

\paragraph{Benchmarks.}
We use three hallucination evaluation benchmarks:
\begin{itemize}
    \item \textbf{FActScore}~\citep{min2023factscore}: Biography generation for 100~entities, with ground-truth atomic facts for verification.
    \item \textbf{LongFact}~\citep{wei2024longform}: Long-form factuality across 50~diverse prompts, with NLI-based labeling using DeBERTa-v3-base-mnli.
    \item \textbf{TruthfulQA}~\citep{lin2022truthfulqa}: 200~questions testing truthfulness in short-form generation.
\end{itemize}

Each model generates responses for all benchmarks, which are decomposed into atomic claims and labeled as factual or hallucinated.
Across models and datasets, we evaluate approximately 460--1,640 claims per model-dataset pair ($\sim$45\% hallucination prevalence on FActScore/LongFact, 50\% on TruthfulQA).

\paragraph{Baselines.}
\begin{itemize}
    \item \textbf{SelfCheckGPT}~\citep{manakul2023selfcheckgpt}: Generate $N{=}5$ alternative responses; check cross-sample consistency via NLI prompting.
    \item \textbf{Verbalized Confidence}: Prompt the model to rate its confidence on a 0--100 scale.
    \item \textbf{Logprob Confidence}: Average token log-probability of the original claim (equivalent to \speccheck{} with $K{=}0$).
    \item \textbf{Random}: Uniform random scores averaged across 3~seeds (42, 123, 456) for stability.
\end{itemize}

Semantic Entropy Probes~\citep{kossen2024semantic}, originally planned as a baseline, were excluded because they require supervised training on annotated data, making direct comparison with our training-free methods inappropriate.

\paragraph{Metrics.}
We report AUC-ROC and AUC-PR for binary hallucination detection.
AUC-PR is more informative when class prevalence varies.
We include 95\% bootstrap confidence intervals computed with 10,000 resamples where indicated.

\subsection{Main Results}

Tables~\ref{tab:main_factscore}--\ref{tab:main_truthfulqa} present per-model results.
Table~\ref{tab:cross_model} summarizes cross-model averages with standard deviations.

\begin{table}[t]
\centering
\caption{Hallucination detection on \textbf{FActScore} (per model). Best in \textbf{bold}. $\uparrow$~=~higher is better.}
\label{tab:main_factscore}
\begin{tabular}{lcccccc}
\toprule
& \multicolumn{2}{c}{Llama-3.1-8B} & \multicolumn{2}{c}{Mistral-7B} & \multicolumn{2}{c}{Qwen2.5-7B} \\
Method & ROC$\uparrow$ & PR$\uparrow$ & ROC$\uparrow$ & PR$\uparrow$ & ROC$\uparrow$ & PR$\uparrow$ \\
\midrule
\speccheck{} & .552 & .488 & .609 & .540 & .623 & .548 \\
SelfCheckGPT & .629 & .573 & \textbf{.733} & \textbf{.660} & \textbf{.723} & \textbf{.662} \\
Verbalized & \textbf{.682} & .581 & .615 & .539 & .707 & .607 \\
Logprob & .697 & \textbf{.585} & .653 & .555 & .689 & .599 \\
Random & .503 & .457 & .508 & .458 & .479 & .436 \\
\bottomrule
\end{tabular}
\end{table}

\begin{table}[t]
\centering
\caption{Hallucination detection on \textbf{LongFact} (per model). Best in \textbf{bold}. $\uparrow$~=~higher is better.}
\label{tab:main_longfact}
\begin{tabular}{lcccccc}
\toprule
& \multicolumn{2}{c}{Llama-3.1-8B} & \multicolumn{2}{c}{Mistral-7B} & \multicolumn{2}{c}{Qwen2.5-7B} \\
Method & ROC$\uparrow$ & PR$\uparrow$ & ROC$\uparrow$ & PR$\uparrow$ & ROC$\uparrow$ & PR$\uparrow$ \\
\midrule
\speccheck{} & .534 & .491 & .514 & .462 & .517 & .455 \\
SelfCheckGPT & .598 & \textbf{.568} & \textbf{.569} & \textbf{.499} & \textbf{.614} & \textbf{.533} \\
Verbalized & .598 & .518 & .547 & .468 & .577 & .488 \\
Logprob & \textbf{.610} & .541 & .526 & .478 & .529 & .464 \\
Random & .517 & .473 & .484 & .449 & .518 & .463 \\
\bottomrule
\end{tabular}
\end{table}

\begin{table}[t]
\centering
\caption{Hallucination detection on \textbf{TruthfulQA} (per model). Best in \textbf{bold}. $\uparrow$~=~higher is better. All methods perform near random chance (prevalence~=~0.50).}
\label{tab:main_truthfulqa}
\begin{tabular}{lcccccc}
\toprule
& \multicolumn{2}{c}{Llama-3.1-8B} & \multicolumn{2}{c}{Mistral-7B} & \multicolumn{2}{c}{Qwen2.5-7B} \\
Method & ROC$\uparrow$ & PR$\uparrow$ & ROC$\uparrow$ & PR$\uparrow$ & ROC$\uparrow$ & PR$\uparrow$ \\
\midrule
\speccheck{} & .486 & .477 & .487 & .489 & .481 & .485 \\
SelfCheckGPT & .510 & \textbf{.515} & \textbf{.539} & .512 & \textbf{.543} & \textbf{.529} \\
Verbalized & \textbf{.521} & .507 & .524 & \textbf{.515} & .540 & .517 \\
Logprob & .503 & .507 & .509 & .508 & .513 & .512 \\
Random & .520 & .515 & .493 & .497 & .504 & .504 \\
\bottomrule
\end{tabular}
\end{table}

\begin{table}[t]
\centering
\caption{Cross-model averages (mean $\pm$ std across 3 models). $\uparrow$~=~higher is better.}
\label{tab:cross_model}
\begin{tabular}{lcccccc}
\toprule
& \multicolumn{2}{c}{FActScore} & \multicolumn{2}{c}{LongFact} & \multicolumn{2}{c}{TruthfulQA} \\
Method & ROC$\uparrow$ & PR$\uparrow$ & ROC$\uparrow$ & PR$\uparrow$ & ROC$\uparrow$ & PR$\uparrow$ \\
\midrule
\speccheck{} & .595{\scriptsize$\pm$.031} & .525{\scriptsize$\pm$.027} & .522{\scriptsize$\pm$.009} & .469{\scriptsize$\pm$.016} & .484{\scriptsize$\pm$.003} & .484{\scriptsize$\pm$.005} \\
SelfCheckGPT & \textbf{.695}{\scriptsize$\pm$.047} & \textbf{.631}{\scriptsize$\pm$.041} & \textbf{.594}{\scriptsize$\pm$.019} & \textbf{.533}{\scriptsize$\pm$.028} & \textbf{.531}{\scriptsize$\pm$.015} & \textbf{.519}{\scriptsize$\pm$.008} \\
Verbalized & .668{\scriptsize$\pm$.039} & .575{\scriptsize$\pm$.028} & .574{\scriptsize$\pm$.021} & .491{\scriptsize$\pm$.021} & .528{\scriptsize$\pm$.008} & .513{\scriptsize$\pm$.005} \\
Logprob & .680{\scriptsize$\pm$.019} & .579{\scriptsize$\pm$.018} & .555{\scriptsize$\pm$.039} & .494{\scriptsize$\pm$.033} & .508{\scriptsize$\pm$.004} & .509{\scriptsize$\pm$.002} \\
Random & .497{\scriptsize$\pm$.013} & .450{\scriptsize$\pm$.010} & .506{\scriptsize$\pm$.016} & .462{\scriptsize$\pm$.010} & .506{\scriptsize$\pm$.011} & .505{\scriptsize$\pm$.007} \\
\bottomrule
\end{tabular}
\end{table}

\paragraph{Key findings.}
(1)~\speccheck{} consistently underperforms SelfCheckGPT, verbalized confidence, and logprob confidence across all model-dataset combinations.
On FActScore, SelfCheckGPT achieves mean AUC-ROC of 0.695 vs.\ 0.595 for \speccheck{}---a 10-point gap.
(2)~On TruthfulQA, all methods perform near random chance (AUC-ROC~$\approx$~0.50), including established baselines.
This suggests that claim-level hallucination detection on TruthfulQA's short-form QA is fundamentally difficult at this model scale.
(3)~\speccheck{}'s best performance is on FActScore with Qwen2.5-7B (AUC-ROC~=~0.623), but even here it falls well short of SelfCheckGPT (0.723).
All \speccheck{} vs.\ SelfCheckGPT bootstrap $p$-values exceed 0.9, confirming \speccheck{} does not outperform SelfCheckGPT on any setting.


%% ============================================================
\section{Analysis}
\label{sec:analysis}
%% ============================================================

\subsection{Why Does Confidence Monotonicity Fail?}
\label{sec:monotonicity_analysis}

Table~\ref{tab:monotonicity} reports the perfect monotonicity rate---the fraction of claims where confidence is strictly non-decreasing across all 4~levels---for factual vs.\ hallucinated claims.

\begin{table}[t]
\centering
\caption{Perfect monotonicity rates (fraction of claims with non-decreasing confidence across all levels) for factual vs.\ hallucinated claims. Higher separation indicates stronger discriminative signal.}
\label{tab:monotonicity}
\begin{tabular}{llccc}
\toprule
Model & Dataset & Factual & Halluc. & Difference \\
\midrule
\multirow{3}{*}{Llama}
  & FActScore  & 0.335 & 0.224 & +0.111 \\
  & LongFact   & 0.462 & 0.382 & +0.080 \\
  & TruthfulQA & 0.377 & 0.385 & $-$0.008 \\
\midrule
\multirow{3}{*}{Mistral}
  & FActScore  & 0.788 & 0.633 & +0.155 \\
  & LongFact   & 0.839 & 0.837 & +0.002 \\
  & TruthfulQA & 0.792 & 0.809 & $-$0.017 \\
\midrule
\multirow{3}{*}{Qwen}
  & FActScore  & 0.797 & 0.622 & +0.175 \\
  & LongFact   & 0.796 & 0.755 & +0.041 \\
  & TruthfulQA & 0.618 & 0.641 & $-$0.023 \\
\bottomrule
\end{tabular}
\end{table}

The results reveal several key issues:

\textbf{Monotonicity holds for most claims regardless of factuality.}
For Mistral and Qwen on FActScore, 63--80\% of \emph{hallucinated} claims exhibit perfect monotonicity.
The hypothesis predicted $<$60\% monotonicity for hallucinated claims and $>$85\% for factual claims; neither threshold is consistently met.

\textbf{Separation is dataset-dependent.}
On FActScore, there is a meaningful gap (11--18 percentage points for Mistral/Qwen), but on LongFact the gap shrinks to 0--8 points, and on TruthfulQA the direction \emph{reverses}: hallucinated claims actually have \emph{higher} monotonicity rates than factual ones.

\textbf{The monotonicity reversal on TruthfulQA} likely occurs because TruthfulQA specifically targets common misconceptions.
The model may be confidently wrong at all specificity levels for these ``common knowledge'' falsehoods, leading to smooth (monotonic) confidence curves even for hallucinated content.

\subsection{Combination with Existing Methods}

Despite its weak standalone performance, \speccheck{} may capture information complementary to existing methods.
Table~\ref{tab:combination} shows results of combining methods via 5-fold cross-validated logistic regression on FActScore (the strongest benchmark for all methods).

\begin{table}[t]
\centering
\caption{Combination results on FActScore via logistic regression (5-fold CV, mean $\pm$ std). Adding \speccheck{} to baselines provides a modest improvement on Mistral and Qwen.}
\label{tab:combination}
\resizebox{\textwidth}{!}{
\begin{tabular}{lcccccc}
\toprule
& \multicolumn{2}{c}{Llama-3.1-8B} & \multicolumn{2}{c}{Mistral-7B} & \multicolumn{2}{c}{Qwen2.5-7B} \\
Combination & ROC$\uparrow$ & PR$\uparrow$ & ROC$\uparrow$ & PR$\uparrow$ & ROC$\uparrow$ & PR$\uparrow$ \\
\midrule
SelfCheck only     & .622{\scriptsize$\pm$.003} & .568{\scriptsize$\pm$.003} & .715{\scriptsize$\pm$.003} & .637{\scriptsize$\pm$.006} & .715{\scriptsize$\pm$.002} & .647{\scriptsize$\pm$.008} \\
All baselines      & .675{\scriptsize$\pm$.004} & .593{\scriptsize$\pm$.003} & .719{\scriptsize$\pm$.003} & .643{\scriptsize$\pm$.004} & .738{\scriptsize$\pm$.001} & .660{\scriptsize$\pm$.004} \\
All + \speccheck{} & .678{\scriptsize$\pm$.002} & .590{\scriptsize$\pm$.002} & \textbf{.740}{\scriptsize$\pm$.003} & \textbf{.659}{\scriptsize$\pm$.002} & \textbf{.747}{\scriptsize$\pm$.001} & \textbf{.681}{\scriptsize$\pm$.001} \\
SC + \speccheck{}  & .649{\scriptsize$\pm$.003} & .571{\scriptsize$\pm$.003} & .735{\scriptsize$\pm$.003} & .654{\scriptsize$\pm$.000} & .727{\scriptsize$\pm$.001} & .664{\scriptsize$\pm$.003} \\
\bottomrule
\end{tabular}}
\end{table}

Adding \speccheck{} to the full baseline set improves AUC-ROC by 0.3 points (Llama), 2.1 points (Mistral), and 0.9 points (Qwen) on FActScore.
The improvement is negligible for Llama but more consistent for Mistral and Qwen, where \speccheck{} adds 1.6--2.1 AUC-PR points.
On LongFact and TruthfulQA (not shown), combination improvements are negligible or negative, consistent with the weak standalone performance on those benchmarks.

\subsection{Ablation: Ladder Depth}

Table~\ref{tab:ablation_depth} shows AUC-ROC for different numbers of specificity levels on FActScore.

\begin{table}[t]
\centering
\caption{Effect of ladder depth $K$ on FActScore AUC-ROC. $K{=}3$ and $K{=}4$ are identical because the $K{=}4$ generation produced the same ladders as $K{=}3$ (see Section~\ref{sec:discussion}).}
\label{tab:ablation_depth}
\begin{tabular}{lcccc}
\toprule
Model & $K{=}1$ & $K{=}2$ & $K{=}3$ & $K{=}4$ \\
\midrule
Llama   & .502 & .567 & .552 & .552 \\
Mistral & .554 & .587 & .609 & .609 \\
Qwen    & .532 & .580 & .623 & .623 \\
\bottomrule
\end{tabular}
\end{table}

Performance generally increases from $K{=}1$ to $K{=}3$ for Mistral and Qwen, with $K{=}2$ performing best for Llama.
$K{=}3$ and $K{=}4$ produce identical results because the model failed to generate meaningful additional specificity levels beyond the 4 already defined, highlighting that the specificity hierarchy has a natural ceiling.

\subsection{Ablation: Score Variants}

Table~\ref{tab:ablation_variants} compares different aggregation strategies for the confidence sequence on FActScore.

\begin{table}[t]
\centering
\caption{Score variant comparison on FActScore (AUC-ROC). The default \speccheck{} score ($\alpha{=}0.5$) performs best for Mistral and Qwen; pure monotonicity ($\alpha{=}0$) is slightly better for Llama.}
\label{tab:ablation_variants}
\begin{tabular}{lccc}
\toprule
Variant & Llama & Mistral & Qwen \\
\midrule
\speccheck{} ($\alpha{=}0.5$) & .552 & \textbf{.609} & \textbf{.623} \\
Pure Monotonicity ($\alpha{=}0$) & \textbf{.564} & .583 & .590 \\
Max Violation     & .544 & .606 & .622 \\
Confidence Gap    & .414 & .427 & .393 \\
\bottomrule
\end{tabular}
\end{table}

The combined \speccheck{} score and pure monotonicity perform similarly, while the confidence gap alone performs \emph{worse than random} (AUC-ROC~$<$~0.5), indicating that the absolute confidence difference between abstract and specific claims does not reliably signal hallucination.

\subsection{Claim Type Analysis}

We categorize claims into entity-specific, numerical, temporal, and general types.
On FActScore, entity-specific claims show the most consistent \speccheck{} signal across models (AUC-ROC 0.58--0.65), while numerical claims show high AUC-PR (0.57--0.71) driven primarily by high hallucination prevalence in that category rather than genuine discriminative power.
Contrary to our hypothesis that \speccheck{} would excel on claims with clear specificity structure (numerical, temporal), the results are inconsistent across models.
Figure~\ref{fig:claim_types} visualizes the per-type breakdown.

\begin{figure}[t]
\centering
\includegraphics[width=0.85\linewidth]{figures/figure_6_claim_types.pdf}
\caption{\speccheck{} AUC-ROC by claim type on FActScore across three models. Entity-specific claims show the most consistent signal. Numerical and temporal claims, where we hypothesized the strongest performance, are inconsistent.}
\label{fig:claim_types}
\end{figure}


%% ============================================================
\section{Discussion and Limitations}
\label{sec:discussion}
%% ============================================================

\paragraph{Why does the hypothesis fail?}
We identify three reasons why confidence monotonicity is a weak hallucination signal:
(1)~\textbf{Confident hallucination}: Models that confidently hallucinate at the specific level tend to remain confident at all abstraction levels, producing monotonic curves for false claims.
(2)~\textbf{Noisy ladder generation}: LLM-generated abstraction introduces semantic drift---the abstract version may not be a strict generalization of the original claim, contaminating the confidence signal.
(3)~\textbf{Logprob saturation}: For claims the model ``wants to say yes to,'' logprob confidence is often near 1.0 at all levels, leaving no room for monotonic increase regardless of factuality.

\paragraph{The N=5/N=10 sampling bug.}
Our ablation comparing sampling-based confidence estimation with $N{=}5$ and $N{=}10$ samples produced identical results across all three random seeds (42, 123, 456).
This indicates a bug in the sampling experiment (likely a seed initialization issue causing both conditions to use the same samples).
We report this transparently rather than presenting flawed ablation results.
The logprob-based method used in all main experiments is unaffected by this issue.

\paragraph{Hallucination Granularity Index.}
Our proposal included a Hallucination Granularity Index to identify \emph{which level of detail} is unreliable.
Given the weak standalone detection performance, the planned 50-case manual evaluation of granularity accuracy was not conducted, as the underlying detection signal is too noisy to support meaningful granularity assessment.
The granularity concept remains theoretically interesting but cannot be validated until the core detection signal improves.

\paragraph{Additional limitations.}
(1)~We evaluate only 7--8B parameter models; larger models may exhibit different confidence patterns.
(2)~FActScore and LongFact labeling relies on automated NLI-based verification, which introduces label noise.
(3)~We use the same model for both generation and ladder construction, which may bias the abstraction toward the model's own patterns.
(4)~Semantic Entropy Probes~\citep{kossen2024semantic} were excluded; future work should compare against trained methods.
(5)~The $K{=}4$ ladder depth was not genuinely different from $K{=}3$, suggesting our abstraction prompting approach has limitations in producing fine-grained intermediate levels.


%% ============================================================
\section{Conclusion}
\label{sec:conclusion}
%% ============================================================

We proposed \speccheck{}, a training-free framework for hallucination detection based on testing confidence monotonicity across specificity levels.
Our extensive experiments across 3~models and 3~benchmarks yield a clear negative result: \speccheck{} underperforms established baselines like SelfCheckGPT (mean AUC-ROC 0.595 vs.\ 0.695 on FActScore, 0.484 vs.\ 0.531 on TruthfulQA).
The core issue is that confidence monotonicity holds for the majority of claims regardless of factuality, providing insufficient discriminative power.

Our analysis reveals that the hypothesis fails because models that hallucinate specific details often do so with consistent confidence across abstraction levels.
However, we find a modest complementary signal when \speccheck{} is combined with existing methods (+0.3--2.1 AUC-ROC points on FActScore), suggesting the specificity dimension captures some independent information.

We believe negative results are valuable for the community: the confidence monotonicity hypothesis is intuitively compelling and might have been independently proposed.
Our thorough evaluation demonstrates its limitations and identifies conditions under which multi-granularity probing may offer marginal value.
Future work could explore whether larger models exhibit stronger monotonicity separation, whether verified abstraction (e.g., using a separate model for ladder generation) yields cleaner signals, or whether the specificity dimension can be better leveraged within ensemble frameworks.


%% ============================================================
\bibliographystyle{plainnat}
\begin{thebibliography}{10}

\bibitem[Farquhar et~al.(2024)]{farquhar2024detecting}
Sebastian Farquhar, Jannik Kossen, Lorenz Kuhn, and Yarin Gal.
\newblock Detecting hallucinations in large language models using semantic entropy.
\newblock \emph{Nature}, 630(8017):625--630, 2024.

\bibitem[Kadavath et~al.(2022)]{kadavath2022language}
Saurav Kadavath, Tom Conerly, Amanda Askell, Tom Henighan, et~al.
\newblock Language models (mostly) know what they know.
\newblock \emph{arXiv preprint arXiv:2207.05221}, 2022.

\bibitem[Kossen et~al.(2024)]{kossen2024semantic}
Jannik Kossen, Jiatong Han, Lorenz Kuhn, Yarin Gal, and Sebastian Farquhar.
\newblock Semantic entropy probes: Robust and cheap hallucination detection in {LLM}s.
\newblock \emph{arXiv preprint arXiv:2406.15927}, 2024.

\bibitem[Lin et~al.(2022)]{lin2022truthfulqa}
Stephanie Lin, Jacob Hilton, and Owain Evans.
\newblock {TruthfulQA}: Measuring how models mimic human falsehoods.
\newblock In \emph{Proceedings of ACL}, 2022.

\bibitem[Manakul et~al.(2023)]{manakul2023selfcheckgpt}
Potsawee Manakul, Adian Liusie, and Mark J.F. Gales.
\newblock {SelfCheckGPT}: Zero-resource black-box hallucination detection for generative large language models.
\newblock In \emph{Proceedings of EMNLP}, 2023.

\bibitem[Min et~al.(2023)]{min2023factscore}
Sewon Min, Kalpesh Krishna, Xinxi Lyu, Mike Lewis, Wen-tau Yih, Pang~Wei Koh, Mohit Iyyer, Luke Zettlemoyer, and Hannaneh Hajishirzi.
\newblock {FActScore}: Fine-grained atomic evaluation of factual precision in long form text generation.
\newblock In \emph{Proceedings of EMNLP}, 2023.

\bibitem[Wang et~al.(2023)]{wang2023selfconsistency}
Xuezhi Wang, Jason Wei, Dale Schuurmans, Quoc~V. Le, Ed~H. Chi, Sharan Narang, Aakanksha Chowdhery, and Denny Zhou.
\newblock Self-consistency improves chain of thought reasoning in language models.
\newblock In \emph{Proceedings of ICLR}, 2023.

\bibitem[Wei et~al.(2024)]{wei2024longform}
Jerry Wei, Chengrun Yang, Xinying Song, Yifeng Lu, Nathan Hu, et~al.
\newblock Long-form factuality in large language models.
\newblock In \emph{Advances in Neural Information Processing Systems (NeurIPS)}, 2024.

\bibitem[Zheng et~al.(2024)]{zheng2024stepback}
Huaixiu~Steven Zheng, Swaroop Mishra, Xinyun Chen, Heng-Tze Cheng, Ed~H. Chi, Quoc~V. Le, and Denny Zhou.
\newblock Take a step back: Evoking reasoning via abstraction in large language models.
\newblock In \emph{Proceedings of ICLR}, 2024.

\end{thebibliography}


%% ============================================================
\appendix
\section{Additional Figures}
\label{sec:appendix}

\begin{figure}[h]
\centering
\includegraphics[width=0.9\linewidth]{figures/figure_1_method_overview.pdf}
\caption{Overview of the \speccheck{} pipeline. An atomic claim is abstracted into a specificity ladder, confidence is estimated at each level via logprob probing, and monotonicity is tested. The top example (factual claim) shows monotonically increasing confidence; the bottom (hallucinated) shows a violation at Level~1.}
\label{fig:method}
\end{figure}

\begin{figure}[h]
\centering
\includegraphics[width=0.9\linewidth]{figures/figure_3_monotonicity_analysis.pdf}
\caption{Monotonicity analysis. Left: distribution of monotonicity scores for factual vs.\ hallucinated claims. Right: average confidence profiles across specificity levels. The substantial overlap between distributions explains the weak discriminative power of the monotonicity signal.}
\label{fig:monotonicity}
\end{figure}

\begin{figure}[h]
\centering
\includegraphics[width=0.9\linewidth]{figures/figure_2_main_results.pdf}
\caption{Main results visualization across all benchmarks. \speccheck{} consistently underperforms SelfCheckGPT and logprob-based confidence.}
\label{fig:main_results}
\end{figure}

%% ============================================================
\section{Reproducibility Details}
\label{sec:reproducibility}
%% ============================================================

\paragraph{Model loading.}
All three models (Llama-3.1-8B-Instruct, Mistral-7B-Instruct-v0.3, Qwen2.5-7B-Instruct) were loaded in FP16 precision (\texttt{torch.float16}) using HuggingFace Transformers with \texttt{device\_map="cuda"} on a single NVIDIA A6000 (48GB).
No quantization was applied.
Left-padding was used for batched inference with decoder-only models.
Random seeds 42, 123, and 456 were used for all stochastic components.

\paragraph{Claim decomposition prompt.}
\begin{quote}\small\ttfamily
Break the following text into independent atomic facts. Each fact should be a single, self-contained sentence that can be verified as true or false independently. Output one fact per line, numbered. Do not include opinions or subjective statements.

Text: \{text\}

Atomic facts:
\end{quote}

\paragraph{Specificity ladder generation prompt (K=3).}
\begin{quote}\small\ttfamily
Given a factual claim, rewrite it at three levels of decreasing specificity. Each level should be less specific but still capture the core meaning.

Example 1:\\
Level 0 (Original): Albert Einstein was born on March 14, 1879 in Ulm, Germany.\\
Level 1 (Approximate): Albert Einstein was born around 1879 in southern Germany.\\
Level 2 (Category): Albert Einstein was born in the late 19th century in Germany.\\
Level 3 (Abstract): Albert Einstein was born in Europe.

[Two additional examples omitted for brevity.]

Now rewrite this claim:\\
Level 0 (Original): \{claim\}\\
Level 1 (Approximate):
\end{quote}
Ladder generation used temperature~0. Invalid ladders (where a more abstract level introduced new specific information) were flagged and re-generated.

\paragraph{Confidence estimation prompt.}
\begin{quote}\small\ttfamily
Is the following statement true? "\{claim\}" Answer with Yes or No.
\end{quote}
Confidence was computed as the softmax probability of ``Yes'' vs.\ ``No'' tokens (summing across casing variants) from a single forward pass.

\paragraph{Verbalized confidence prompt.}
\begin{quote}\small\ttfamily
Rate your confidence that the following statement is true on a scale of 0 to 100, where 0 means certainly false and 100 means certainly true. Only output the number.\\
Statement: "\{claim\}"
\end{quote}

\paragraph{Generation parameters.}
Biography and long-form generation used \texttt{max\_new\_tokens=256} (FActScore) or \texttt{512} (LongFact), with \texttt{temperature=0.7}, \texttt{top\_p=0.9}.
TruthfulQA answers used \texttt{max\_new\_tokens=128}.
SelfCheckGPT generated $N{=}5$ alternative responses per prompt at \texttt{temperature=1.0}, \texttt{top\_p=0.95}.

\paragraph{Computational cost.}
The full pipeline for one model takes approximately 3--4 hours on a single A6000: generation and claim decomposition ($\sim$50 min), specificity ladder generation ($\sim$20 min), confidence estimation ($\sim$10 min per dataset), and baseline computation ($\sim$60 min for SelfCheckGPT, $\sim$5 min each for verbalized and logprob baselines).
Total wall-clock time for all 3 models, all experiments, and ablations was approximately 14 hours.


\end{document}
